% Unidad U010. Espejo orientado y selección 3->1; redacción autónoma.

\chapter{Mirror involution and local selection \(3\to1\) of the surviving extension}
\label{chap:espejo-seleccion-tres-uno}

The ninth phase contains two local input states. Once the oriented
state \(B_9\) has been fixed, the boundary relation produces three extensions
toward the first phase. The three share an axis of self-scale and a ternary
aggregate; they differ in the ordered distribution of that aggregate between the branches of
closure and propagation. A second horizon conserves exactly one.

\section{Symmetric and axial coordinates}

Let
\[
 \mathcal B:=(\mathbb F_3^6)^3.
\]
A visible block is
\[
 B=(u^\pi,u^e,u^\varphi)\in\mathcal B.
\]
We define
\[
 q(B):=u^\pi+u^e+u^\varphi,
 \qquad
 a(B):=u^\pi+u^e-u^\varphi.
\]
All operations are performed componentwise in
\(\mathbb F_3^6\).

\begin{theorem}[Reconstruction of the axis and the aggregate]
\label{thm:espejo-reconstruccion-eje-agregado}
The coordinates \(q,a\) determine
\[
 \boxed{u^\varphi=2(q-a),}
 \qquad
 \boxed{u^\pi+u^e=q-u^\varphi.}
\]
\end{theorem}

\begin{proof}
Subtracting the definitions gives
\(q-a=2u^\varphi\). In \(\mathbb F_3\), \(2^{-1}=2\); therefore,
\(u^\varphi=2(q-a)\). Substitution into the definition of \(q\) recovers the
remaining aggregate.
\end{proof}

The specular involution is
\[
 \mathfrak c(u^\pi,u^e,u^\varphi)
 =(u^e,u^\pi,u^\varphi).
\]

\begin{proposition}[Invariants of the Involution]
\label{prop:espejo-invariantes}
\(\mathfrak c^2=\operatorname{id}\), y
\[
 q(\mathfrak cB)=q(B),
 \qquad
 a(\mathfrak cB)=a(B).
\]
In particular, \(u^\varphi\) and \(u^\pi+u^e\) are invariantes.
\end{proposition}

\begin{proof}
Exchanging the first coordinates twice recovers \(B\). Both
\(q\) and \(a\) depend on them only through their sum, so that the
exchange preserves them. The \cref{thm:espejo-reconstruccion-eje-agregado}
transfers that invariance to the axis and the aggregate.
\end{proof}

The involution asserts neither that \(\pi=e\) nor that one is the quotient of the other.
It asserts that the two branches occupy conjugate positions with respect to a
fixed structure.

\section{Oriented state in phase \(9\) and associated boundary data}

The visible projection of the selected state is
\[
 B_9=(100100\mid020112\mid010122).
\]
The induced relation on visible blocks is denoted
\[
 \mathcal M_9^B\subseteq\mathcal B\times\mathcal B.
\]
The two local input states of the phase table precede this
relation. After selecting the complete state that projects to \(B_9\), its
output fiber has cardinal three.

\begin{theorem}[The three extensions of the ninth phase]
\label{thm:espejo-tres-extensiones}
The relational image of \(B_9\) contains exactly
\begin{align*}
B_{10}^{(0)}&=(101222\mid112221\mid112002),\\
B_{10}^{(1)}&=(102221\mid111222\mid112002),\\
B_{10}^{(2)}&=(111221\mid102222\mid112002).
\end{align*}
The three satisfy
\[
 q_{10}=022112,
 \qquad
 a_{10}=101111,
\]
and, consequently,
\[
 u_{10}^\varphi=112002,
 \qquad
 u_{10}^\pi+u_{10}^e=210110.
\]
\end{theorem}

\begin{proof}
The extension relation on complete states is enumerated in the fiber of
\(B_9\) and produces the three indicated rows. Adding and subtracting their three
components gives the same pair \((q_{10},a_{10})\) in all three cases. The
\cref{thm:espejo-reconstruccion-eje-agregado} produces the common axis and
aggregate. The first two components are distinct, so the three
rows represent different distributions of the same aggregate.
\end{proof}

All their publications in base three are
\[
 (279,352,365),
 \qquad
 (280,351,365),
 \qquad
 (282,350,365).
\]

\section{Exact compatibility of cylinders}

A ternary word of length \(L\) and index \(N\) determines the half-open
cylinder
\[
 I_3(N,L)=\left[\frac{N}{3^L},\frac{N+1}{3^L}\right).
\]
A prefix of \(b\) decimal blocks, with associated integer \(D\), determines
\[
 I_{1000}(D,b)=
 \left[\frac{D}{1000^b},\frac{D+1}{1000^b}\right).
\]

\begin{lemma}[Integer intersection test]
\label{lem:espejo-test-entero-cilindros}
The cylinders \(I_3(N,L)\) and \(I_{1000}(D,b)\) intersect if and only if
\[
 N1000^b<(D+1)3^L,
 \qquad
 (N+1)1000^b>D3^L.
\]
\end{lemma}

\begin{proof}
Two half-open intervals \([a,b)\) and \([c,d)\) intersect if and only if \(a<d\) and \(c<b\). We substitute the four endpoints and multiply by the positive integer \(3^L1000^b\).
\end{proof}

The test uses only integer arithmetic. It does not depend on a floating-point tolerance.

\section{2nd horizon and uniqueness}

The following block is
\[
 B_{11}=(022212\mid110001\mid100021),
\]
and the corresponding decimal boundary is
\[
 (502,662,638).
\]

\begin{theorem}[Local resolution \(3\to1\)]
\label{thm:espejo-seleccion-tres-uno}
The three extensions of \cref{thm:espejo-tres-extensiones} are compatible at depth one. When
\(B_{11}\) and the boundary \((502,662,638)\) are imposed simultaneously,
only \(B_{10}^{(0)}\) prolongs as the same complete state. Therefore,
\[
 \boxed{
 \#\operatorname{Surv}^{\rm cal}_1(B_9)=3,
 \qquad
 \#\operatorname{Surv}^{\rm cal}_2(B_9)=1.}
\]
\end{theorem}

\begin{proof}
We apply \cref{lem:espejo-test-entero-cilindros} to each of the three
channels. At the first horizon, the six inequalities of every candidate are
satisfied. After adjoining \(B_{11}\), the inequalities of the three channels
remain satisfied for \(B_{10}^{(0)}\). In \(B_{10}^{(1)}\) and
\(B_{10}^{(2)}\), compatibility of channel \(\varphi\) is preserved, but
the conditions of branches \(\pi\) and \(e\) fail. Enumeration therefore
leaves a unique state at depth two.
\end{proof}

This proof does not define the global connection from a local example. It is an
explicit witness of how a multivalued relation at one horizon can
become single-valued by preserving memory and requiring a longer prolongation.

\section{Phase return with block change}

The first phase of the initial loop is
\[
 B_1=(012222\mid121221\mid112202),
\]
while the first phase of the next turn is
\[
 B_{10}=(101222\mid112221\mid112002).
\]
Thus,
\[
 g(1)=g(10)=1,
 \qquad
 B_1\neq B_{10}.
\]
The phase returns, but the axis, the aggregate, the distribution, and the memory of
quotient have evolved.

\section{Diagram of the mirror involution \(\pi/e\), the fixed axis \(\varphi\), and the local selection \(3\to1\) by prolongability}

\begin{figure}[htbp]
\centering
\resizebox{0.98\textwidth}{!}{%
\begin{tikzpicture}[
 node distance=10mm and 15mm,
 every node/.style={font=\small,align=center},
 state/.style={draw=black,rounded corners,inner sep=5pt},
 inv/.style={draw=black,double,rounded corners,inner sep=5pt},
 arr/.style={-{Latex[length=2mm]},semithick}
]
 \node[state] (pi) {$u^\pi$\\closure branch};
 \node[state,right=70mm of pi] (e) {$u^e$\\propagation branch};
 \coordinate (mirroraxis) at ($(pi)!0.5!(e)$);
 \node[inv,above=9mm of mirroraxis] (phi) {$u^\varphi$\\eje fijo};
 \draw[arr,<->] (pi) -- node[below] {$\mathfrak c$} (e);
 \node[above=8mm of phi] (aggregate)
   {$u^\pi+u^e$ preserved in the fiber};
 \draw[dashed] (phi.north) -- (aggregate.south);
 \node[state,below=20mm of mirroraxis] (b9) {phase $9$: $B_9$\\two local inputs};
 \node[state,below left=18mm and 25mm of b9] (b100) {$B_{10}^{(0)}$};
 \node[state,below=18mm of b9] (b101) {$B_{10}^{(1)}$};
 \node[state,below right=18mm and 25mm of b9] (b102) {$B_{10}^{(2)}$};
 \node[state,below=18mm of b101] (b11) {$B_{11}$\\unique prolongation};
 \draw[arr] (b9) -- (b100);
 \draw[arr] (b9) -- (b101);
 \draw[arr] (b9) -- (b102);
 \draw[arr] (b100) -- (b11);
 \draw[arr,densely dotted] (b101) -- node[right] {obstructed} (b11);
\draw[arr,densely dotted] (b102) -- node[right] {obstructed} (b11);
\end{tikzpicture}
}
\caption{The mirror preserves the axis \(u^\varphi\) and the aggregate
\(u^\pi+u^e\). The return of the ninth phase opens three visible extensions;
the memory of the second horizon preserves one. The circuit does not reset the
state.}
\label{fig:espejo-tres-uno-arrastre}
\end{figure}

\section{Residue--quotient memory during the mirror}

The recurrence
\[
 x_kA=u_k+3x_{k+1}
\]
of the \cref{eq:bockstein-hensel-local} is applied to each of the three
channels. The part \(u_k\) publishes the visible block; \(x_{k+1}\) conserves
the quotient. Since \(A\in\operatorname{GL}_6(\mathbb Z)\), the pair recovers
the preceding state. The mirror can exchange the visible branches without
eliminating their provenance quotients.

\section{Obstructions to the mirror involution \(\pi/e\) and to the local selection \(3\to1\): invariants, depth-\(2\) horizon, and Hensel memory}

\begin{criterion}[Refutation Criteria for the Mirror]
The construction is refuted if:
\begin{enumerate}
 \item the involution modifies \(u^\varphi\) or \(u^\pi+u^e\);
 \item the two local input states are conflated with the three outputs;
 \item one of the three outputs has a different \((q_{10},a_{10})\);
 \item the cylinder test uses a floating tolerance;
 \item more than one output survives the declared second horizon;
 \item phase return is interpreted as equality \(B_1=B_{10}\);
 \item the visible exchange erases the Hensel quotients.
\end{enumerate}
\end{criterion}

The mirror and the local selection have now been constructed. The following
units will develop the Sturmian vacancies and then the three
complete analytic biographies.
